\documentclass[10pt,a4paper]{amsart}
\usepackage[margin=19mm]{geometry}
\usepackage{amsmath,amssymb,mathtools}
\usepackage[scr=rsfs,cal=euler]{mathalfa}
\usepackage{xfrac}
\usepackage[unicode,hidelinks,bookmarks=false]{hyperref}
\newcommand{\ie}{i.e.\ }
\newcommand{\cf}{cf.\ }
\newcommand{\resp}{respectively\ }
\newcommand{\Subsection}{\subsection}
\providecommand{\nobreakdash}{\nobreak\hbox{-}}
\setlength{\emergencystretch}{2em}
\multlinegap0pt
\usepackage{bm}
\usepackage{bbm}
\usepackage{dsfont}
\usepackage{xspace}
\usepackage{cancel}
\usepackage{mathtools}
\usepackage{amsthm}
\makeatletter
\@ifundefined{theorem}{\newtheorem{theorem}{Theorem}}{}
\@ifundefined{corollary}{\newtheorem{corollary}[theorem]{Corollary}}{}
\@ifundefined{proposition}{\newtheorem{proposition}[theorem]{Proposition}}{}
\makeatother
\newcommand\quadcover[1]{\mathcal{U}^\square_{#1}}
\title[Short-range and long-range order: a transition in block-glui]{Short-range and long-range order: a transition in block-gluing behavior in Hom shifts}
\author{Silvere Gangloff, Benjamin Hellouin de Menibus, Piotr Oprocha}
\date{Source: arXiv:2211.04075v3 (CC BY 4.0)}
\begin{document}
\maketitle

\noindent\emph{Source and licence.} Short-range and long-range order: a transition in block-gluing behavior in Hom shifts. Version arXiv:2211.04075v3 (CC BY 4.0), distributed under the Creative Commons Attribution 4.0 International licence, as recorded in the arXiv licence metadata for this exact version and shown on the abstract page https://arxiv.org/abs/2211.04075v3. Full rights notice: licenses/arxiv-2211.04075-CC-BY-4.0.txt.\par\smallskip

\noindent\textit{Editorial scope.} Counted unit: the Proposition whose source label is \texttt{proposition.square.cover.decomposable} in the article (source0.tex lines 1053--1055, source label \texttt{proposition.square.cover.decomposable}), delivered together with the article's own complete original proof of it (source0.tex lines 1057--1095, one \texttt{proof} environment of 39 lines). No mathematics is written, repaired or re-derived: statement and proof are the article's own text line for line. Required context retained uncounted: cor.morphism.quotient (lines 849--851). Every definition, display and label of those lines is retained unchanged. Omitted from this package and not redistributed: 1--848, 852--1052, 1056--1056, 1096--2478. Nothing inside a retained excerpt is omitted or altered: every retained body is delivered exactly as the article's lines are written, so no omission receipt of any kind accompanies this package. Citations: the retained excerpts carry no citation command at all, so no bibliography entry of the article is transcribed and no reference is redistributed. Reference adaptation: none was needed and none was made; every reference inside the delivered unit resolves inside it, so no unresolved reference or citation is produced. Printed slips kept literal: no printed word, symbol, hypothesis or number of the counted statement or of its proof was altered. Layout and class: the article's own class is \texttt{article} with packages including inputenc, authblk, latexsym, amsthm, amscd, amsmath, amssymb, amsfonts, tikz, stmaryrd, hyperref, tcolorbox, enumerate, units; the wrapper is the lane's pinned \texttt{amsart} editorial wrapper at 10pt on A4, which restates the theorem-like environments the retained text needs and adds the article's own macro definitions that the retained text needs (\texttt{\textbackslash quadcover}), with extra wrapper packages (bm, bbm, dsfont, xspace, cancel, mathtools). The delivered numbering is the unit's own. Assets: no retained span contains a figure environment, a graphics-inclusion command, a PDF-page inclusion, a table or a TikZ picture; no image file is copied into this package, the package declares no asset dependency and no image file is redistributed with it. Licence: version arXiv:2211.04075v3 of this article is distributed under the Creative Commons Attribution 4.0 International licence, as recorded in the arXiv licence metadata for this exact version and shown on the abstract page https://arxiv.org/abs/2211.04075v3. A full rights notice is delivered at licenses/arxiv-2211.04075-CC-BY-4.0.txt. Characters: every retained excerpt is pure ASCII, so the delivered engine, XeLaTeX with its default Latin Modern text font, reports no missing character.
\begin{corollary}\label{cor.morphism.quotient}
For every $a, b$ neighbors in $G$ and $p,q \in V_{\mathcal{U}_G[b]}$ such that $p \sim_{\square} q$, we have that $\psi_{b \mapsto a}(p) \sim_{\square} \psi_{b \mapsto a}(q)$.
\end{corollary}

\begin{proposition}\label{proposition.square.cover.decomposable}
The square cover of $G$ is square-decomposable.
\end{proposition}

\begin{proof}
%Let us label $\quadcover{G}$ by walks starting from some vertex $a$ of $G$, that is, $\quadcover{G}[a]$.
%Let us consider some vertex $a$ of $G$.
We will use the following: if $c$ is a cycle on $\quadcover{G}$, then $\eta(c)$ is a cycle, and $c$ has a backtrack iff $\eta(c)$ has a backtrack as well.

%Remember that, for any vertex $v$ of $\quadcover{G}[a]$ (which is an equivalence class for $\sim_{\square}$), $\eta(v)$ is the final vertex of all walks in $v$. We generalize the notation to walks $w$ on $\quadcover{G}[a]$, so that $\eta(w)$ is a walk on $G$. 

Take any vertex $v$ of $\quadcover{G}$. By repeated application of morphisms $\overline{\psi}$, we can choose a labeling of $\quadcover{G}$ so that $v$ is the class for $\sim_\square$ of the empty walk.
Let $w$ be a walk on $G$ such that $w_0 = \eta(v)$. We denote by $\eta^{-1}_v(w)$ the walk such that for all $i \le l(w)$:
\[\eta^{-1}_v(w)_i\text{ is the class for }\sim_\square\text{ of the walk } \varphi(w_0\dots w_i).\]
It is straightforward that for any walk $w'$ on $\quadcover{G}$, $\eta^{-1}_{w'_0}\circ\eta(w') = w'$.

%For any vertex $v$ of $\quadcover{G}$, let $u$ be a non-backtracking walk on $G$ such that $v = \pi_a(u)$ and %let $w$ be a walk on $G$ such that $w_0 = \eta(v)$, we denote by $\eta^{-1}_v(w)$ the walk such that for all %$i \le l(w)$:
%\[\eta^{-1}_v(w)_i\text{ is the class for }\sim_\square\text{ of the walk } \varphi(u\odot w_1\dots w_i).\]
%It is straightforward that for any walk $w'$ on $\quadcover{G}[a]$, $\eta^{-1}_{w'_0}\circ\eta(w') = w'$.

%\eta \circ \eta^{-1}_v(w) = w$ for any $w$ $v$ such that  $w_0 = \eta(v)$.

Let us turn to the proof. Consider a cycle $c$ on $\quadcover{G}$. We will prove that it is square-decomposable.
\begin{enumerate}
\item \textbf{The cycle $\eta(c)$ is square-decomposable:}

Since $c$ is a non-backtracking cycle, $\eta(c)$ is also a non-backtracking cycle. 
Furthermore, since $c_0 = c_{l(c)}$, there exists $p$ in $c_0$ such that $\varphi(p \eta(c)) \in c_0$, therefore 
$\varphi(p \eta(c)) \sim_{\square} p$.
Then we can use Corollary~\ref{cor.morphism.quotient} repeatedly to have 
\[\psi_{p_{l(p)-1} \mapsto p_{l(p)}} \circ \cdots \circ \psi_{p_0 \mapsto p_1} (\varphi(p \eta(c))) \sim_{\square} \psi_{p_{l(p)-1} \mapsto p_{l(p)}} \circ \cdots \circ \psi_{p_0 \mapsto p_1} (p).\]

This is rewritten into $\eta(c) \sim_{\square} p_{l(p)}$ (empty cycle), which means that $\eta(c)$ is square-decomposable. Let us consider $(q^{(i)})_{0\le i\le l}$ a square decomposition of $\eta(c)$.

\item \textbf{The sequence $(\eta^{-1}_{c_0}(q^{(i)}))_{0\le i\le l}$ is a square decomposition of $c$:} 

It is clear that $\eta^{-1}_{c_0}(\eta(c)) = c$ and $\eta^{-1}_{c_0}(\eta(c_0)) = c_0$ (empty cycle). Therefore, it is sufficient to prove that for all $i$ such that $0\leq i < l$, $\eta^{-1}_{c_0}(q^{(i)})$ and $\eta^{-1}_{c_0}(q^{(i+1)})$ differ by a square.

Considering such an integer $i$, we can assume that $q^{(i+1)} = q^{(i)} \oplus_k s$ for some square $s$ and index $k$ (the other possible case is processed similarly). For any vertex $v$ of $\quadcover{G}$ such that $\eta(v) = s_0$, we have that $\eta^{-1}_{v}(s)$ is a square in $\quadcover{G}$. Indeed, it is non-backtracking, and it is a cycle since $v$ and $v\odot s$ differ by a square.

It is straightforward to check that $\eta^{-1}_{c_0}(q^{(i+1)}) = \eta^{-1}_{c_0}(q^{(i)}) \oplus_k \eta^{-1}_{v}(s)$, where $v =  \eta^{-1}_{c_0}(q^{(i)})_k$. Since $\eta^{-1}_{v}(s)$ is a square, we have that $\eta^{-1}_{c_0}(q^{(i)})$ and $\eta^{-1}_{c_0}(q^{(i+1)})$ differ by a square. This concludes the proof.\qedhere
\end{enumerate}
\end{proof}

\end{document}
